\section{Parameters Studied}
\label{sec:parameters}

\subsection{METIS Imbalance Factor ($\ufactor$)}

The METIS graph partitioner~\citep{karypis2013manual} accepts a parameter
\texttt{ufactor} that controls the maximum allowable imbalance in the
partition: each partition is required to have weight between
$(1 - \ufactor) \cdot P_{\rm target}$ and $(1 + \ufactor) \cdot P_{\rm target}$,
where $P_{\rm target} = P_{\rm total} / k$.

In the redistricting context, METIS is applied at each node of the
\ar{} bisection tree.
The DIA population tolerance of $\pm 0.5\%$ is a separate constraint
applied at the leaf level (individual district).

The $\ufactor$ range $[0.1\%, 5\%]$ covers:
\begin{itemize}
\item $0.1\%$: near-perfect balance within METIS; maximally constrains
      the partitioner, reducing its ability to find compact cuts.
\item $1\%$ (DIA default): balances compactness and balance.
\item $5\%$: wide METIS tolerance; allows very compact cuts at the cost
      of greater population imbalance within METIS subtrees (the DIA
      leaf-level constraint still applies).
\end{itemize}

\subsection{County Edge Weight ($\acounty$)}

County-boundary edges receive a multiplicative weight $\acounty$ relative
to standard geographic boundary-length edges.
This implements the ``county-respecting'' principle from B.10~\citep{b10paper}:
the algorithm prefers cuts along county lines.

The range $\acounty \in [1.0, 10.0]$:
\begin{itemize}
\item $1.0$: county edges are weighted the same as non-county edges;
      county-respecting behavior is disabled.
\item $3.0$ (DIA default): county lines are preferred 3$\times$ over
      equally long non-county edges.
\item $10.0$: strong county-respecting; the algorithm will almost always
      cut along county lines, even at some compactness cost.
\end{itemize}

\subsection{ConvergenceSweep Threshold ($T$)}

$T$ is the number of consecutive non-improving seeds required to halt
the sweep.
The range $T \in [100, 1{,}000]$:
\begin{itemize}
\item $T = 100$: fast but may miss the global minimum (B.16 shows
      Georgia last improves at seed 489, so $T = 100$ fails GA).
\item $T = 600$ (DIA default): certifiably sufficient for all 50 states
      with $P(\text{miss}) < 0.001$.
\item $T = 1{,}000$: extra margin; no additional benefit for any state.
\end{itemize}

\subsection{Area Swing ($\aswing$)}

In GeoSection/AreaSection (B.8/B.9), the area swing parameter controls
the relative weight of geographic area versus population in the bisection
objective.
The range $\aswing \in [1.05, 1.20]$:
\begin{itemize}
\item $1.05$: near-population-proportional bisection.
\item $1.10$ (DIA default): moderate area weight.
\item $1.20$: strong area weight; the algorithm will bisect more
      spatially than demographically.
\end{itemize}

\subsection{VRA Boost Weight ($\wvra$)}

VRA-eligible tracts (census blocks with $\geq 40\%$ minority population)
receive edge weight boost $\wvra$ to reduce the probability of splitting
minority communities.
The range $\wvra \in [0.2, 0.6]$:
\begin{itemize}
\item $0.2$: weak VRA boost; minority community protection is minimal.
\item $0.4$ (DIA default): moderate VRA boost.
\item $0.6$: strong VRA boost; significant compactness cost.
\end{itemize}
